\section{Overview}
\label{sec:overview}
Before presenting the full details of our type system, we begin with a
tour of its key features and how they enable both safety and
reachability reasoning.

%% Readers can see this
%% connection via the type erasure function $\eraserf{\tau}$, which
%% converts context types to basic types by erasing these refinements. In
%% the following paragraphs, we introduce these refinements in detail via
%% examples.

\paragraph{Demonic and Angelic Modalities}
% BD: The goal of this subsection is to equip the reader with the
% right intuitions of what the modalities mean when ascribed to types
% and when used in the typing context of a typing judgment. I'm
% currently of the opinion that the lower / upper bound framing is the
% right intuition (particularly for the type context case), as opposed
% to the angelic / demonic one.

% The reader should have also internalized the motto'that "angelic
% refinement = reachability, demonic refinement = safety" by the end
% of this subsection.

% Note that fundamental question that our type system answers is: what
% does the image of this term look like, when it is evaluated under a
% particular restriction? One of the interesting things about the type
% system is a) the sorts of restrictions our type contexts can
% describe, and b) how those restrictions afford the type certain
% reasoning capabilities when verifying safety / reachability
% properties.

As in other refinement type systems, a context type refines a base
type $b$ with a qualifier $\phi$, a sentence in first-order
logic. Every context type is equipped with a modality indicating
whether it should be interpreted either demonically ($\ort{b}{\phi}$)
or angelically ($\urt{b}{\phi}$). When ascribed to a term $e$, a
context type describes both safety and reachability properties of
$e$. As discussed in \autoref{sec:intro}, the demonic context type
$\ort{b}{\phi}$ allows $e$ to (at most) acquire any value $v$ that
satisfies its qualifier $\phi$.  Conversely, an angelic context type
$\urt{b}{\phi}$ stipulates that $e$ must (at least) acquire every
value satisfying $\phi$.

% Informally, we have
% \begin{align*}
%   &\vdash e : \ort{b}{\phi} \quad\text{iff}\quad \forall v. \mstep{e}{v} \implies \phi[\vnu \mapsto v]
%   \\&\vdash e : \urt{b}{\phi} \quad\text{iff}\quad \forall v. \phi[\vnu \mapsto v] \implies  \mstep{e}{v}
% \end{align*}\noin

To illustrate this difference, consider the following typing
judgements which ascribe different context types to the constant term
$1$:
\begin{align*}
  &\vdash 1 : \ort{\Nat}{\top}\;\judgok
  && \vdash 1 : \urt{\Nat}{\top}\;\judgfail
  \\&\vdash 1 : \ort{\Nat}{\vnu > 0}\;\judgok
  && \vdash 1 : \urt{\Nat}{\vnu > 0}\;\judgfail
  \\&\vdash 1 : \ort{\Nat}{\vnu = 1}\;\judgok
  && \vdash 1 : \urt{\Nat}{\vnu = 1}\;\judgok
\end{align*}\noindent
The judgements on the left assert \emph{safety} properties, similar to
a standard refinement type system, where the constant $1$ needs to
satisfy the corresponding qualifier. The judgements on the right, in
contrast, describe \emph{reachability} constraints: since the constant
$1$ cannot reach every natural number or positive number, the first
two judgements are invalid. The last (singleton) type precisely
captures the reachability properties the term exhibits. Under an empty
typing context, a deterministic term %
% \footnote{All terms in \langname{} are deterministic, a restriction
% that we will lift in \autoref{sec:extension}.}
can only be assigned either an empty or singleton angelic type, as it
can evaluate to at most one value.  Even for deterministic terms,
however, richer reachability properties naturally arise when reasoning
within a broader program context. To do so, our type system
\emph{holistically} considers the \emph{set} of values that a term can
evaluate to, i.e., its image, under a restricted \emph{set} of
environments, or \emph{capability}, that is constrained by a typing
context.
% Put another way, our type system describes the image of a term under
% a set of environments consistent with a typing context.

Intuitively, a demonic refinement acts as an \emph{upper} bound: when
used as a result type, $\ort{b}{\phi}$ requires every value in the
image of that term to satisfy $\phi$; in a type context, it only
admits capabilities in which every environment maps the variable to a
value satisfying $\phi$. An angelic refinement, in contrast, acts as a
\emph{lower} bound: as a result type, $\urt{b}{\phi}$ requires the
image of a term to contain every value satisfying $\phi$; in a type
context, it forces capabilities to at least include environments
mapping the variable to each value that satisfies $\phi$. To
illustrate, % these
% bounds manifest in the guarantees provided by our type system,
consider the following typing judgements involving non-singleton
reachability properties: %
\begin{align*}
  x{:}\urt{\Nat}{\vnu > 3} \vdash x - 3 : \urt{\Nat}{\vnu >
  0}\;\judgok \quad & \quad
  x{:}\urt{\Nat}{\vnu > 4} \vdash x - 3 : \urt{\Nat}{\vnu > 0}\;\judgfail
\end{align*} %
\noindent
Clearly, for any particular positive number $n$, the term $x - 3$ will
evaluate to that $n$, assuming $x$ can be assigned $n + 3$.  The
angelic types of $x$ in each typing context place different
constraints on the set of possible assignments to $x$ in a capability:
under the first, for every $n$ greater than $3$, the capability must
contain an environment in which $x$ is $n$; the second allows a
capability which does not include a mapping of $x$ to $4$, and thus
cannot guarantee the image of $x - 3$ includes $1$.

Switching the modality on the type of $x$ inverts the direction of
these bounds, thereby restricting the set of assignments to $x$ that
are impossible. As a consequence, the reachability guarantees that our
type system can provide about deterministic terms involving only
demonically qualified variables are somewhat limited:
\begin{align*}
  x{:}\ort{\Nat}{\vnu > 2} \vdash x - 3 : \urt{\Nat}{\vnu > 0}\;\judgfail % \\
  % x{:}\ort{\Nat}{\vnu > 4} \vdash x - 3 : \urt{\Nat}{\vnu > 0}\;\judgfail
\end{align*} %
Observe there are many singleton capabilities containing an
environment in which $x$ is assigned a value greater than $2$, but
$x - 3$ can only reach a single result value under such a capability!
% If all we know about $x$ is that it cannot be less than $2$, after
% all, there's little we can say about the values that $x-3$ can reach!
Similarly, our type system can provide somewhat limited safety
guarantees when only equipped with reachability assumptions:
\begin{align*}
  x{:}\urt{\Nat}{\vnu > 4} \vdash x - 3 : \ort{\Nat}{\vnu > 0}\;\judgfail
\end{align*} %
The image of $x - 3$ includes a non-positive value under any
capability containing an environment where $x$ is $4$, and the lower
bound provided by the angelic type of $x$ in the typing judgement
above does not preclude such an environment. Using our type system, it
is still possible to derive meaningful safety guarantees from angelic
assumptions through the use of the $\oplus$ operator, as we shall see
shortly.

%% \BD{At some point early on, we need to introduce the notion of method
%%   predicates. This may just happen organically with the addition of
%%   the buggy merge example to the introduction... }

% This
% A typing
% judgement can ascribe a non-trivial (i.e., non-singleton) angelic type
% to its term based on its typing context.  For example, the following
% type judgement is valid:
% \begin{align*}
%   x{:}\urt{\Nat}{\top} \vdash x - 3 : \urt{\Nat}{\vnu > 0}\;\judgok
% \end{align*}\noindent
% Clearly, the term $x - 3$, when evaluated, can produce all positive
% numbers because its corresponding type context $x{:}\urt{\Nat}{\top}$
% assumes that assignments to variable $x$ can reach all natural
% numbers; in particular, this means $x$ can be any number greater than
% 3.  An angelic type context allows the type system to presume a
% desirable assignment of $x$ depending on the shape of the
% right-hand-side type; e.g., in order to reach $n$, the assignment of
% $x$ should be $n + 3$.

% We assume an \emph{external impure environment} that serves as a source
% of angelic inputs.  For example, a test generator that yields random
% numbers can serve as an example of such a source:
% \begin{align*}
%   \Envir_x(e) \doteq \zlet{x}{\Code{nat\_gen()}}{\boxed{e}}
% \end{align*}
% For our purposes here, we do not include environments such as
% $\Envir_x$ as part of our core language.

\paragraph{Subtyping}
% \BD{TBH, the material in this subsection does not seem
%   particularly novel (the subtyping relations are similar to prior
%   work) or important to what follows (I don't believe swapping
%   modalities on singleton refinements is not used elsewhere). }

The subtyping relationship between two angelic or two demonic
refinements is similar to other contemporaneous refinement type
systems~\cite{LiquidHaskell, CoverageType}, where the subtype
relationship between demonic refinements is covariant with the
qualifier, and that of angelic refinements is contravariant with the
qualifier. These relationships align with the interpretation of
angelic and demonic modalities as lower and upper bounds: a term
which can reach all positive numbers can also reach all numbers
greater than 1, while a term that always safely evaluates to a number
greater than 1 is also guaranteed to always evaluate to a positive
number.
\begin{align*}
  &\ort{\Nat}{\vnu > 1} \subty \ort{\Nat}{\vnu > 0}\;\judgok
  &&\urt{\Nat}{\vnu > 1} \subty \urt{\Nat}{\vnu > 0} \;\judgfail \\
  & \ort{\Nat}{\vnu > 0} \subty \ort{\Nat}{\vnu > 1}  \;\judgfail
  && \urt{\Nat}{\vnu > 0} \subty \urt{\Nat}{\vnu > 1}\;\judgok
\end{align*} %

The subtyping relationship between angelic and demonic refinement
types is more nuanced. No subtyping relationship exists between two
arbitrary refinements with different modalities: %
\begin{align*}
  &\ort{\Nat}{\vnu > 0} \subty \urt{\Nat}{\vnu > 0}\;\judgfail
  & &\urt{\Nat}{\vnu > 0} \subty \ort{\Nat}{\vnu > 0} \;\judgfail
\end{align*} %
\noindent
It is only possible to switch the modality of a refinement when its
qualifier is a singleton, or when the resulting refinement type
trivially holds, i.e., $\ort{b}{\top}$ or $\urt{b}{\bot}$: %
\begin{align*}
  &\ort{\Nat}{\vnu = 3} \subty \urt{\Nat}{\vnu =
    3}\;\judgok
  & & \urt{\Nat}{\vnu = 3} \subty \ort{\Nat}{\vnu =
      3}\;\judgfail \\
  & \urt{\Nat}{\vnu > 0} \subty \ort{\Nat}{\top}\;\judgok
  & &\ort{\Nat}{\vnu > 0} \subty \urt{\Nat}{\bot}\;\judgok
\end{align*}\noindent
When the qualifier is a singleton, safety entails reachability ---
there is only one possible value, so the demonic constraint (the value
may only be 3) and the angelic constraint (the value 3 must be
reachable) make identical demands on the term. Converting to a
supertype with a trivial qualifier, in contrast, simply throws away
the bounds on a term entirely.
%% \SJ{I'm not sure why the following paragraphs are related to subtyping -
%%   the explanations don't explicitly refer to the subtyping relation, and
%%   the may/must example seems too complicated at this stage.}
%% Knowing that a term can only evaluate to positive numbers does not
%% allow us to draw any conclusions about which positive numbers it
%% reaches, for example. Similarly, simply knowing that a term can reach
%% every positive number does not preclude it from also evaluating to 0.
%% \begin{wrapfigure}{r}{0.48\linewidth}
%%   \centering
%%   \begin{minted}[linenos = false, fontsize = \footnotesize, escapeinside=??]{OCaml}
%% val g : ?$\Code{i}{:}\ort{\Int}{\vnu \le 10}\ctsarr\ort{\Int}{\vnu \neq 0}$?

%% ?$\vdash$? let f (x: int) (y: int) =
%%   if y > 10 then (g x) + 1
%%     else x + 1
%% : ?$\Code{x}{:}\urt{\Int}{\vnu \le 5}\ctsarr\Code{y}{:}\urt{\Int}{\vnu \le 15}\ctsarr\urt{\Int}{\vnu = 1}$?
%%   \end{minted}
%%   \caption{A simple example of how may-analysis helps
%%     must-analysis~\cite{CompositionalMayMust}.}
%%   \label{fig:may-must-reach}
%% \end{wrapfigure}
%% That said, numerous prior works~\cite{GHK+06,BVR+08,GN+10,NYC+12} have shown that knowledge
%% about the safety properties of terms can aid the automated
%% verification of reachability properties, and vice-versa. When
%% establishing that the function \Code{f} in
%% \autoref{fig:may-must-reach} must return \Code{1} on some arguments,
%% for example, the program analysis of \citet{GN+10} uses
%% a not-may summary analogous to the signature of \Code{g} to conclude
%% calling it will never produce the \Code{0} value needed for \Code{f}
%% to yield \Code{1}. As we shall see shortly, our type system also
%% naturally supports using angelic information in the type context to
%% verify a meaningful form of qualified safety properties, similar to
%% that found in \citet{ZDS23} and \citet{RHLE}.

% that an angelic context requiring $x$ to reach all numbers greater
% than 1 is a subtype of one requiring $x$ to reach all positive
% numbers. This holds because any term that is typed under the
% assumption that $x$ can reach every number greater than 1 can
% certainly be given the same type under the assumption that $x$ can
% reach every positive number, since the latter subsumes the former.
% In other words, a narrower angelic qualifier (fewer reachable
% values) is a weaker assumption, while a wider angelic qualifier
% (more reachable values) is stronger — the opposite of the demonic
% case.

% Having both angelic as well as demonic modalities in a type context
% requires a more refined subtype relation than usual.  Read
% ``$\Gamma_1 \ctsubseteqA \Gamma_2$'' as ``any term that has type
% $\tau$ under $\Gamma_1$ can also be typed with $\tau$ under $\Gamma_2$
% .''  Then, the following statements hold:
% \begin{align*}
%   &x{:}\ort{\Nat}{\top} \ctsubseteqA x{:}\ort{\Nat}{\vnu > 0}\;\judgok
%   &&x{:}\urt{\Nat}{\top} \ctsubseteqA x{:}\urt{\Nat}{\vnu > 0} \;\judgfail
%   \\&x{:}\ort{\Nat}{\vnu > 1} \ctsubseteqA  x{:}\ort{\Nat}{\vnu > 0}  \;\judgfail
%   &&x{:}\urt{\Nat}{\vnu > 1} \ctsubseteqA x{:}\urt{\Nat}{\vnu > 0}\;\judgok
% \end{align*}\noindent
% A demonic (over-approximate) context has a subtype relation that is
% covariant with the qualifier, while the angelic (under-approximate)
% context is contravariant with the qualifier.  For example, the last
% statement says that an angelic context requiring $x$ to reach all
% numbers greater than 1 is a subtype of one requiring $x$ to reach all
% positive numbers. This holds because any term that is typed under the
% assumption that $x$ can reach every number greater than 1 can
% certainly be given the same type under the assumption that $x$ can
% reach every positive number, since the latter subsumes the former.  In
% other words, a narrower angelic qualifier (fewer reachable values) is
% a weaker assumption, while a wider angelic qualifier (more reachable
% values) is stronger — the opposite of the demonic case.

% Additionally, the subtype relation must allow context shifts between
% these two modalities since they serve as resources used by the type
% system.  Specifically, a demonic context cannot be treated as a
% non-trivial angelic context:
% \begin{align*}
%   &x{:}\ort{\Nat}{\vnu > 0} \ctsubseteqA x{:}\urt{\Nat}{\top}\;\judgfail
%   &&x{:}\ort{\Nat}{\vnu > 0} \ctsubseteqA x{:}\urt{\Nat}{\vnu > 1}\;\judgfail
%   \\&x{:}\ort{\Nat}{\vnu = 3} \ctsubseteqA x{:}\urt{\Nat}{\vnu = 3}\;\judgok
%   &&x{:}\ort{\Nat}{\vnu > 0} \ctsubseteqA x{:}\urt{\Nat}{\bot}\;\judgok
% \end{align*}\noindent
% For example, the demonic context $x{:}\ort{\Nat}{\vnu > 0}$ only
% constrains $x$ to be \emph{some} positive number, providing no
% guarantee about which positive numbers are reachable.  Treating it as
% a non-trivial angelic context would require it to serve as a witness
% that \emph{every} value satisfying the angelic qualifier is reachable,
% a guarantee that a demonic context is fundamentally incapable of
% providing.   Demonic and angelic contexts align when
% their type qualifier is trivially false ($\bot$) or is a singleton
% (e.g., $\vnu = 3$). When the qualifier is a singleton, safety and
% reachability coincide — there is only one possible value, so the
% demonic constraint (the value must be 3) and the angelic constraint
% (the value 3 must be reachable) make identical demands on the term.
% Finally, an angelic context can be viewed as a subtype
% of a demonic one provided every value reachable under the angelic
% qualifier also satisfies the demonic qualifier.
% \begin{align*}
%   x{:}\urt{\Nat}{\vnu > 0} \ctsubseteqA x{:}\ort{\Nat}{\top}\;\judgok
%   &&x{:}\urt{\Nat}{\vnu > 0} \ctsubseteqA x{:}\ort{\Nat}{\vnu > 1}\;\judgfail
% \end{align*}\noindent
% The statement on the right is incorrect because the angelic context
% guarantees that $x$ can reach every positive number, including
% $1$. However, the demonic context $x{:}\ort{\Nat}{\vnu > 1}$ only
% permits values strictly greater than 1 so a term typed as the
% singleton $1$ under the angelic context would violate the demonic
% constraint.

\paragraph{Type Context Composition}
% Here we want to transition from thinking of type contexts as
% constraining the surrounding environment to thinking of them as
% affording us certain reasoning capabilities. Importantly, if we want
% to reason compositionally, we need to have a finer-grained notion of
% the sorts of capabilities it can provide
When they appear in a typing context, the bounds imposed by our two
modalities afford us different reasoning capabilities: when
establishing the reachability properties of a term, for example, an
angelic refinement lets us \emph{choose} a favorable assignment for a
variable in that term, while a demonic refinement allows us to
\emph{rule out} certain assignments to that variable when showing that
a term is safe.
% BD: If there's room, we could point back to specific examples in the
% previous subsection to emphasize this point.
Importantly, the bounds imposed by the bindings of individual
variables in a typing context can interact with each other, altering
the capabilities afforded by the typing context as a whole. To
illustrate, consider the following pair of program contexts:
\begin{align}
  &\zlet{x}{\Code{gen_{nat} () }}{\zlet{y}{x + 1}{\square}}
    \tag{C$_\mathit{en}$}\label{ex:EnCon} \\
  &\zlet{x}{\Code{gen_{nat} () }}{\zlet{y}{\Code{gen_{nat}()}}{\square}}
    \tag{C$_\mathit{dis}$}\label{ex:DisjCon}
\end{align}\noindent
Both contexts use $\Code{gen_{nat} ()}$ to (nondeterministically)
generate an arbitrary value for $x$. The first context then sets $y$
based on the value of $x$, while the second opts to generate a fresh
value for $y$ instead. The values bound to $x$ and $y$ in either
context can all reach every possible positive number, as can the
program \ref{ex:DisjCon}[$x - y$].  This is not true of
\ref{ex:EnCon}[$x - y$], however, because the values of $x$ and $y$ in
this program are \emph{entangled} since the values $y$ can yield are
tied to $x$'s.
\begin{align*}
  &\emptyset \vdash \text{\ref{ex:DisjCon}}[x] : \urt{\Nat}{\vnu > 0}
    \;\judgok \quad
  &\emptyset \vdash \text{\ref{ex:EnCon}}[x] : \urt{\Nat}{\vnu > 0}
    \;\judgok  \\
  & \emptyset \vdash \text{\ref{ex:DisjCon}}[y] : \urt{\Nat}{\vnu > 0}
  \;\judgok \quad
  &\emptyset \vdash \text{\ref{ex:EnCon}}[y] : \urt{\Nat}{\vnu > 0}
    \;\judgok \\
  & \emptyset \vdash \text{\ref{ex:DisjCon}}[x - y] : \urt{\Nat}{\vnu
    > 0} \;\judgok \quad & \emptyset \vdash \text{\ref{ex:EnCon}}[x - y] : \urt{\Nat}{\vnu > 0} \;\judgfail
\end{align*}\noindent
To deal with this situation, our type system includes two constructors
over typing contexts, ``$\ctcomma$'' and ``$\ctstar$'', each of which
affords different capabilities when reasoning about, e.g., subtyping
relationships. Informally, $\Gamma_1 \ctstar \Gamma_2$ means that the
capabilities of $\Gamma_1$ and $\Gamma_2$ are independent; when either
context includes an angelic type binding, we are free to ``choose'' a
favorable assignment for the corresponding variable independent of the
other context.  Dually, $\Gamma_1 \ctcomma \Gamma_2$ means that the
capabilities of $\Gamma_1$ and $\Gamma_2$ are dependent; the angelic
choices afforded by $\Gamma_2$ may be constrained by choices for
variables in $\Gamma_1$.  Consider the following two subtyping
judgments reflecting the typing of $\square $ in $C_{\mathit{dis}}$
and $C_{\mathit{en}}$, resp:
\begin{align*}
  &x{:}\urt{\Nat}{\vnu \geq 0} \ctstar y{:}\urt{\Nat}{\vnu \geq 1}
    \vdash \urt{\Nat}{\vnu = x - y} \subty \urt{\Nat}{\vnu > 0}
    \;\judgok \tag{$\Gamma_\mathit{dis}$}\label{eqn:Ctxdis} \\
  &x{:}\urt{\Nat}{\vnu \geq 0} \ctcomma y{:}\urt{\Nat}{\vnu \geq 1}
    \vdash \urt{\Nat}{\vnu = x - y} \subty \urt{\Nat}{\vnu > 0} \;\judgfail
    \tag{$\Gamma_\mathit{en}$}\label{eqn:Ctxen}
\end{align*}\noindent
Under \ref{eqn:Ctxdis}, we can freely choose both $2$ for $y$
\emph{and} an arbitrary value larger than $2$ for $x$, from which the
subtyping relationship follows. Under \ref{eqn:Ctxen}, however, the
possible values we can choose for $x$ and $y$ are constrained by the
\emph{entangled} conjunction, ``$\ctcomma$'': once a value is chosen
for $x$, $y$ can no longer be freely chosen. Concretely, the angelic
guarantee for $y$ under ``$\ctcomma$'' only promises that $y$ can
reach every value satisfying $\vnu > 1$ \emph{for some fixed
  instantiation of $x$}.  Since the subtyping goal requires producing
every positive number as $x - y$, and this requires pairing each
candidate value of $x$ with an independently chosen $y$, a dependent
context cannot be used to supply such pairs.  This is evidenced by
$C_{en}$, which always collapses its result to 0, rather than all
positive numbers.  The ``\ctcomma'' in the context
$\Gamma\ctcomma \ x{:}\tau_x \ctcomma y{:}\tau_y$ works like additive
logic conjunction in BI, and both
$\Gamma\ctcomma \ x{:}\tau_x \ctcomma y{:}\tau_y \vdash x{:}\tau_x$
and
$\Gamma\ctcomma \ x{:}\tau_x \ctcomma y{:}\tau_y \vdash y{:}\tau_y$
are valid. Importantly, however, it does not guarantee that the
modalities of $x$ and $y$ hold \emph{at the same
  time}. \footnote{Unlike a linear type system, which would
  \emph{consume} the capability to angelically assign $x$ when
  producing the capability to angelically assign $y$, our system
  instead \emph{entangles} these capabilities, allowing the body
  expression $e$ to use both $x$ and $y$ while still tracking their
  dependency. This formulation aligns with the intuition of bunched
  typing~\cite{BunchedType}. % \BD{I don't think the reader is going to
    % appreciate this difference unless they have a better intuition of
    % the capabilities afforded by $\ctcomma$.}
}  For example,
\begin{align*}
  &x{:}\urt{\Nat}{\vnu > 0} \ctstar (y{:}\urt{\Nat}{\vnu > 0} \ctcomma z{:}\urt{\Nat}{\vnu > 0}) \vdash x - y : \urt{\Nat}{\vnu > 0} \;\judgok
  \\&x{:}\urt{\Nat}{\vnu > 0} \ctstar (y{:}\urt{\Nat}{\vnu > 0}
  \ctcomma z{:}\urt{\Nat}{\vnu > 0}) \vdash x - (z + y) : \urt{\Nat}{\vnu > 0} \;\judgok
  \\&x{:}\urt{\Nat}{\vnu > 0} \ctstar (y{:}\urt{\Nat}{\vnu > 0} \ctcomma z{:}\urt{\Nat}{\vnu > 0}) \vdash y - z : \urt{\Nat}{\vnu > 0} \;\judgfail
\end{align*}\noindent
In these judgements, the context tracks variable dependency,
specifying that $x$ is independent from $y$ and $z$, while $y$ and $z$
are dependent on each other.  Specifying independence in this way is
sufficient to verify reachability for $x - y$ and $x - (z + y)$, since
$x$ can be angelically chosen independently of both $y$ and $z$.
However, type-checking $y - z$ against the type $\urt{\Nat}{\vnu > 0}$
fails because $y$ and $z$ are entangled: no single angelic choice can
simultaneously make their difference reach every positive number.

Notably, angelic and demonic refinements can also be entangled in a
way that impacts the capabilities afforded by a type context, as the
following pair of judgments illustrate:
\begin{align*}
  &x{:}\urt{\Nat}{\vnu \geq 0} \ctcomma y{:}\ort{\Nat}{\vnu \geq 1} \vdash x - y : \urt{\Nat}{\I{even}(\vnu)} \;\judgfail
  \\&x{:}\urt{\Nat}{\vnu \geq 0} \ctstar y{:}\ort{\Nat}{\vnu \geq 1} \vdash x - y : \urt{\Nat}{\I{even}(\vnu)} \;\judgok
\end{align*}\noindent
The first typing context stipulates that the value of $y$ is positive,
but that value is allowed to depend on the value $x$ --- the hole in
\ref{ex:EnCon} can also be typed under this typing
context. Accordingly, the expression \ref{ex:EnCon}$[x - y]$ cannot
reach every even number. Contrast this with a variant of this context
that instead assigns \emph{some} constant value $n$ to $y$:
\begin{align}
  &\zlet{x}{\Code{gen_{nat} () }}{\zlet{y}{n}{\square}}
    \tag{C$_{n}$}\label{ex:DisjSafeCon}
\end{align}
Regardless of the particular value of $n$, for any given even number
$m$, choosing $x = m + n$ allows \ref{ex:DisjSafeCon}$[x - y]$ to
reach that number.

% \BD{The $\ctsubseteqA$ relation is doing some heavy lifting in our
%   type system, but at this point, it is not clear how that relation
%   interacts with additive and multiplicative conjunction. Maybe we can
%   provide some algebraic laws to help build up the reader's
%   intuitions?}


\paragraph{Function Types}
Entangled ($\ctsarr$) and independent
function types ($\ctwand$) are adjunctions of the corresponding type
context composition operators:
\begin{align*}
  \Gamma \vdash e_f : x{:}\tau_x \ctsarr \tau \quad &\text{ iff } \quad \Gamma \ctcomma x{:}\tau_x \vdash (e_f\;x) : \tau
  \\\Gamma \vdash e_f : x{:}\tau_x \ctwand \tau \quad &\text{ iff } \quad \Gamma \ctstar x{:}\tau_x \vdash (e_f\;x) : \tau
\end{align*}\noindent
Thus, ($\Code{-}$) can be given the following types:
\begin{align*}
  &\vdash (-) : x{:}\urt{\Nat}{\vnu > 0} \ctwand y{:}\urt{\Nat}{\vnu > 0} \ctwand \urt{\Nat}{\vnu > 0}
  \\&\vdash (-) : x{:}\urt{\Nat}{\vnu > 0} \ctsarr y{:}\urt{\Nat}{\vnu > 0} \ctwand \urt{\Nat}{\vnu > 0}
\end{align*}\noindent
Notably, it cannot have the following type:
\begin{align*}
  &\vdash (-) : x{:}\urt{\Nat}{\vnu > 0} \ctwand y{:}\urt{\Nat}{\vnu > 0} \ctsarr \urt{\Nat}{\vnu > 0}
\end{align*}\noindent
Doing so would require proving the following type judgement:
\begin{align*}
  &\emptyset \ctstar x{:}\urt{\Nat}{\vnu > 0} \ctcomma y{:}\urt{\Nat}{\vnu > 0} \vdash x - y : \urt{\Nat}{\vnu > 0} \;\judgfail
\end{align*}\noindent
As we saw earlier, a type context in which $x$ and $y$ are entangled
restricts the angelic choices that can be made, preventing us from
establishing that $x - y$ can produce every positive number.

Higher-order functions introduce additional complexity. Consider the
following example:
\begin{align*}
  &x{:}\urt{\Nat}{\vnu > 0} \ctstar y{:}\urt{\Nat}{\vnu > 0} \vdash \zzlam{f}{x - (f\; (f\;y))} :
  \\&\quad  \underbrace{f{:}(\underbrace{y{:}\urt{\Nat}{\vnu > 0} \ctsarr}_\texttt{$y$ can depend on context during application} \urt{\Nat}{\vnu > 0}) \ctwand}_\texttt{$f$ is independent from $x$ and $y$} \urt{\Nat}{\vnu > 0}
\end{align*}\noindent
To make $x - (f\;(f\;y))$ reach every positive number, $x$ and
$f\;(f\;y)$ must be independent. To ensure this, one might expect that
all function arrows in the result type would need to be $\ctwand$, but
this would be incorrect. If $f$ had the fully independent type
$y{:}\urt{\Nat}{\vnu > 0} \ctwand \urt{\Nat}{\vnu > 0}$, then each
application of $f$ would require its argument to be independent from
$f$'s closure. While the first application $f\;y$ satisfies this, the
second application $f\;(f\;y)$ fails because the intermediate result
$f\;y$ depends on $f$. While $\ctwand$ guarantees that $f$'s argument
is not entangled with its closure, it makes no guarantees that $f$'s
result is also independent. Thus, while $(f\; y)$ is independent of
$y$, it may still depend on $f$'s closure, making it an invalid
argument in the outer application when equipped with this signature.

If $f$ is instead given the entangled function type
$y{:}\urt{\Nat}{\vnu > 0} \ctsarr \urt{\Nat}{\vnu > 0}$, its argument
can now depend on $f$'s closure, accommodating the nested
application. The function itself is then typed with $\ctwand$,
ensuring $f$'s body is not entangled with either $x$ or $y$ in the
outer context. By adjunction, this yields:
\begin{align*}
  &x{:}\urt{\Nat}{\vnu > 0} \ctstar
  \\&\quad (y{:}\urt{\Nat}{\vnu > 0} \ctstar
  f{:}(y{:}\urt{\Nat}{\vnu > 0} \ctsarr \urt{\Nat}{\vnu > 0}))
  \vdash x - (f\; (f\;y)) : \urt{\Nat}{\vnu > 0}
\end{align*}\noindent
Since $x$ is independent from both $y$ and $f$, it is also independent
from $f\;(f\;y)$, thus guaranteeing that $x - (f\;(f\;y))$ can reach
every positive number.

These examples illustrate a flexibility that, to our knowledge, is
absent from prior refinement type systems: demonic and angelic
qualifiers may appear in arbitrary combination throughout a type, and
the entangled ($\ctsarr$) and independent ($\ctwand$) function type
constructors may likewise be chosen independently at each argument
position.

\paragraph{From Reachability to Safety}
As discussed above, the sort of safety that can be directly derived
from reachability assumptions in the typing context is quite limited
--- ascribing a demonic refinement to a term constrains \emph{all} of
its possible executions, while an angelic type binding only makes
assumptions about a \emph{subset} of the environments that a term may
be run in. It is possible, however, to derive meaningful safety
properties about the executions that an angelic refinement guarantees
exist. If we have assumed that $x$ must reach every even number, for
example, we know that the term $x \Code{+ 1}$ will always safely
evaluate to an odd number \emph{under those environments}. In order to
soundly establish upper bounds on program executions that are
guaranteed to exist, our type system includes a \emph{context sum}
type constructor $\ctoplus$ that explicitly identifies partitions of
the values a term must reach.

The typing judgement $\Gamma \vdash e : \tau_1 \ctoplus \tau_2$ states
that $e$ can be typed with $\tau_1$ and $\tau_2$ in some partitioning
of $\Gamma$, and that all capabilities consistent with $\Gamma$ are
accounted for by $\tau_1$ and $\tau_2$ --- no possibility is lost. For
example: \begin{align*}
  &x{:}\urt{\Nat}{\vnu < 10} \vdash x + 1 : \ort{\Nat}{odd(\vnu)} \ctoplus
    \ort{\Nat}{even(\vnu)} \;\judgok \notag \\
  &x{:}\urt{\Nat}{\top} \vdash \zif{x == 0}{3}{5} : \ort{\Nat}{\vnu <
    4}\ctoplus \ort{\Nat}{\vnu > 4} \;\judgok
    %\tag{\Code{e}$_\oplus$}\label{ex:oplus}
\end{align*}\noindent
In both judgements, the result type uses $\ctoplus$ to explicitly
demarcate two subsets of the term's executions. The first example uses
$\ctoplus$ to establish that the angelic refinement of $x$ in the
typing context ensures that $x + 1$ must evaluate to at least one even
and at least one odd number, without stipulating what those numbers
are. In the second judgment, $\ctoplus$ is used to distinguish between
the two possible control flows through the program.

Importantly, while
$\ort{\Nat}{\vnu < 4} \ctoplus \ort{\Nat}{\vnu > 4}$ is a subtype of
$\ort{\Nat}{\vnu < 4 \lor \vnu > 4}$ (both types can only be ascribed
to a term that never evaluates to $4$), the reverse is not true: the
disjunctive refinement type loses the information that a term has two
distinct feasible control flow paths, one whose result is always
greater than $4$ and another whose result is always less than
$4$.\footnote{In this
  sense, the $\ctoplus$  operator plays a role analogous to the commutative
  monoid disjunction in Outcome Logic~\cite{ZDS23}.} This ability also
allows us to focus on the safety properties of one particular feasible
program path, assigning all remaining paths the $\top$ type: %
\begin{align*}
  &x{:}\urt{\Nat}{\top} \vdash \zif{x > 0}{5}{3} : \ort{\Nat}{\vnu >
    4} \ctoplus \ort{\Nat}{\top} \quad \judgok
    \tag{\Code{e}$_{\oplus\top}$}\label{ex:oplusWeak}
\end{align*}\noindent
The use of $\tau \ctoplus \ort{b}{\top}$ expresses an \emph{affine}
style of control flow analysis when deriving information from angelic
facts: like affine logic, which allows resources to be used at most
once (and hence discarded), this style allows the type system to
discard uninteresting branches by weakening them to $\ort{b}{\top}$,
imposing a precise safety constraint on the branches that a typing
context guarantees exist while placing no constraint on the others.

% Standard refinement types interpret type
% contexts in a demonic style, requiring \emph{all} control-flow within
% the program to be consistent with the safety property.  In the
% presence of angelic contexts that can reach multiple control flows,
% additional handling is required.  Consider the following judgements:
% \begin{align*}
%   &x{:}\urt{\Nat}{\top} \vdash \zif{x == 0}{3}{5} : \ort{\Nat}{\vnu < 4} \quad \text{(reach true branch)}
%   \\&x{:}\urt{\Nat}{\top} \vdash \zif{x == 0}{3}{5} : \ort{\Nat}{\vnu > 4} \quad \text{(reach false branch)}
% \end{align*}\noindent
% These judgements informally state that ``there exists an input natural
% number $x$ such that we may reach some number less than (greater than,
% resp.) $4$.''  They mix safety and reachability reasoning, encoding a
% reachability property using imprecise overapproximated constraints
% (i.e., $\ort{\Nat}{\vnu < 4}$ and $\ort{\Nat}{\vnu > 4}$) instead of
% precise underapproximated constraints (i.e., $\urt{\Nat}{\vnu = 3}$
% and $\urt{\Nat}{\vnu = 5}$), since the desired safety property does not need
% to be precise about the exact result values.

% However, na\"{i}vely combining these imprecise reachability properties
% with safety reasoning may lead to semantic conflicts.  A purely safety-focussed
% interpretation of the first judgement would require that
% \emph{every} value produced by $\zif{x == 0}{3}{5}$ under the
% angelic context $x{:}\urt{\Nat}{\top}$ satisfies $\vnu < 4$.  This is
% incorrect: since $x$ can angelically take any natural number,
% including $1$, the expression can evaluate to $5$, violating $\vnu < 4$.
% This conflict also breaks the standard semantic characterization
% of intersection types found in refinement type systems that are used to model
% conditionals:
% \begin{align*}
%   &\Gamma \vdash e : \ort{b}{\phi_1} \ctinter \ort{b}{\phi_2} \text{ iff } \Gamma \vdash e : \ort{b}{\phi_1 \land \phi_2}
% \end{align*}\noindent
% since no term can have type $\ort{\Nat}{\vnu < 4 \land \vnu > 4}$.
% The root cause is that an angelic context cannot be treated as simple
% existential quantification (e.g., ``there exists an $x$ that reaches
% the true branch'') that discards the possibility that $x$ can also
% reach the false branch.  In order to preserve all possibilities while
% still enabling an existential-like interpretation of angelic contexts,
% we introduce a \emph{context sum} type constructor $\ctoplus$ that
% explicitly tracks  expected control flows.
